% ==============================================================
% Section 2: Epstein-Zin Preferences
% ==============================================================

\section{Epstein-Zin Preferences}
\label{sec:epstein-zin}

The standard time-separable expected utility model ties together two conceptually distinct aspects of preferences: aversion to risk (across states) and aversion to intertemporal fluctuations (across time). Epstein-Zin preferences untie this knot, allowing risk aversion and the elasticity of intertemporal substitution to be specified independently. This separation is not merely a generalization for its own sake---it resolves empirical puzzles and provides a cleaner conceptual framework for dynamic portfolio choice.

% --------------------------------------------------------------
\subsection{The Recursive Aggregator}
\label{subsec:ez-aggregator}
% --------------------------------------------------------------

\paragraph{Motivation.}

Consider two distinct questions an investor might face:
\begin{enumerate}[nosep]
    \item \textbf{Risk aversion}: How much would you pay to avoid a fair gamble over consumption \emph{today}?
    \item \textbf{Intertemporal substitution}: How willing are you to shift consumption between \emph{today and tomorrow} in response to interest rate changes?
\end{enumerate}

These are logically independent questions. An investor might be highly risk averse (unwilling to bear uncertainty) yet also highly willing to substitute intertemporally (responsive to interest rates). Or vice versa.

With standard time-separable CRRA utility,
\begin{equation}
U = \E_0 \left[ \sum_{t=0}^{\infty} \beta^t \frac{c_t^{1-\gamma}}{1-\gamma} \right],
\label{eq:time-separable}
\end{equation}
both questions are answered by the \emph{same} parameter $\gamma$: risk aversion equals $\gamma$, and the elasticity of intertemporal substitution (IES) equals $1/\gamma$. This is a strong restriction with no economic justification.

\paragraph{The Epstein-Zin Aggregator.}

Epstein and Zin (1989) proposed a recursive utility specification that separates these two roles. The value function $V_t$ satisfies
\begin{equation}
\boxed{V_t = \left[ (1-\beta) \, c_t^{1-1/\theta} + \beta \left( \E_t\left[ V_{t+1}^{1-\gamma} \right] \right)^{\frac{1-1/\theta}{1-\gamma}} \right]^{\frac{1}{1-1/\theta}},}
\label{eq:epstein-zin}
\end{equation}
where:
\begin{itemize}[nosep]
    \item $\theta > 0$ is the \textbf{elasticity of intertemporal substitution} (IES)
    \item $\gamma > 0$ is the \textbf{coefficient of relative risk aversion}
    \item $\beta \in (0,1)$ is the \textbf{discount factor}
\end{itemize}

The recursion has two nested aggregations, each with CES structure.

\paragraph{Two Nested CES Aggregations.}

Equation \eqref{eq:epstein-zin} combines two generalized means:

\begin{enumerate}[label=(\roman*)]
    \item \textbf{Aggregation across states} (inner): The term
    \begin{equation}
    \mu_t(V_{t+1}) \equiv \left( \E_t\left[ V_{t+1}^{1-\gamma} \right] \right)^{\frac{1}{1-\gamma}}
    \label{eq:certainty-equivalent-v}
    \end{equation}
    is the \emph{certainty equivalent} of next period's value. This is a generalized mean with power $\rho_{\text{risk}} = 1 - \gamma$ and weights equal to probabilities. Risk aversion $\gamma$ governs this aggregation.
    
    \item \textbf{Aggregation across time} (outer): Given the certainty equivalent $\mu_t$, the recursion
    \begin{equation}
    V_t = \left[ (1-\beta) \, c_t^{1-1/\theta} + \beta \, \mu_t^{1-1/\theta} \right]^{\frac{1}{1-1/\theta}}
    \label{eq:intertemporal-aggregation}
    \end{equation}
    is a CES aggregator over $(c_t, \mu_t)$ with power $\rho_{\text{time}} = 1 - 1/\theta$. The IES $\theta$ governs this aggregation.
\end{enumerate}

\begin{calloutimportant}[The Two-Layer Structure]
Epstein-Zin preferences nest two CES aggregations:
\begin{center}
\begin{tabular}{@{}lcc@{}}
\toprule
\textbf{Aggregation} & \textbf{Power Parameter} & \textbf{Elasticity} \\
\midrule
Across states (risk) & $\rho_{\text{risk}} = 1 - \gamma$ & $1/\gamma$ \\
Across time (intertemporal) & $\rho_{\text{time}} = 1 - 1/\theta$ & $\theta$ \\
\bottomrule
\end{tabular}
\end{center}
The certainty equivalent $\mu_t$ serves as the ``price'' of future utility in the intertemporal aggregation. This nested structure is the key to tractability.
\end{calloutimportant}

\paragraph{Rewriting with Explicit Means.}

Using our notation for generalized means, \eqref{eq:epstein-zin} becomes
\begin{equation}
V_t = \M\left( c_t, \mu_t(V_{t+1}); (1-\beta, \beta), 1-\tfrac{1}{\theta} \right),
\label{eq:ez-as-mean}
\end{equation}
where
\begin{equation}
\mu_t(V_{t+1}) = \Mbar\left( V_{t+1}; \mathbf{p}_t, 1-\gamma \right)
\label{eq:mu-as-mean}
\end{equation}
and $\mathbf{p}_t$ is the vector of state probabilities conditional on information at $t$.

% --------------------------------------------------------------
\subsection{The Time-Separable Case}
\label{subsec:time-separable}
% --------------------------------------------------------------

\paragraph{The Restriction $\gamma = 1/\theta$.}

A remarkable simplification occurs when risk aversion equals the inverse of the IES:
\begin{equation}
\gamma = \frac{1}{\theta} \quad \Longleftrightarrow \quad \rho_{\text{risk}} = \rho_{\text{time}} = 1 - \gamma.
\label{eq:ez-restriction}
\end{equation}
Under this restriction, both aggregations use the \emph{same} power parameter. The nested structure collapses.

\begin{proposition}[Recovery of Time-Separable Utility]
\label{prop:time-separable}
When $\gamma = 1/\theta$, Epstein-Zin utility reduces to standard time-separable CRRA expected utility:
\begin{equation}
V_t = \left( \E_t \left[ \sum_{s=t}^{\infty} \beta^{s-t} c_s^{1-\gamma} \right] \right)^{\frac{1}{1-\gamma}}.
\label{eq:ez-to-crra}
\end{equation}
Equivalently, maximizing $V_t$ is equivalent to maximizing $\E_t[\sum_{s=t}^{\infty} \beta^{s-t} u(c_s)]$ with $u(c) = c^{1-\gamma}/(1-\gamma)$.
\end{proposition}

\begin{proof}
With $\gamma = 1/\theta$, write $\rho = 1 - \gamma = 1 - 1/\theta$. The recursion \eqref{eq:epstein-zin} becomes
\begin{align}
V_t &= \left[ (1-\beta) c_t^{\rho} + \beta \left( \E_t[V_{t+1}^{\rho}] \right)^{\rho/\rho} \right]^{1/\rho} \notag \\
&= \left[ (1-\beta) c_t^{\rho} + \beta \, \E_t[V_{t+1}^{\rho}] \right]^{1/\rho}.
\label{eq:collapsed-recursion}
\end{align}
Define $\tilde{V}_t = V_t^{\rho}$. Then
\[
\tilde{V}_t = (1-\beta) c_t^{\rho} + \beta \, \E_t[\tilde{V}_{t+1}].
\]
This is a linear recursion with solution
\[
\tilde{V}_t = (1-\beta) \E_t \left[ \sum_{s=t}^{\infty} \beta^{s-t} c_s^{\rho} \right].
\]
Taking the $1/\rho$ power and noting that monotonic transformations preserve optima yields \eqref{eq:ez-to-crra}.
\end{proof}

\paragraph{Why This Restriction is Problematic.}

The restriction $\gamma = 1/\theta$ implies:
\begin{itemize}[nosep]
    \item High risk aversion ($\gamma > 1$) $\Rightarrow$ low IES ($\theta < 1$): Risk-averse investors are also unwilling to substitute intertemporally.
    \item Low risk aversion ($\gamma < 1$) $\Rightarrow$ high IES ($\theta > 1$): Risk-tolerant investors are also highly responsive to interest rates.
\end{itemize}

Empirically, these implications are problematic:
\begin{itemize}[nosep]
    \item The equity premium puzzle suggests $\gamma \approx 10$ or higher to rationalize observed risk premia.
    \item But $\theta = 1/\gamma = 0.1$ implies extremely low intertemporal substitution, inconsistent with consumption responses to interest rate changes.
\end{itemize}

Epstein-Zin preferences resolve this by allowing, say, $\gamma = 10$ and $\theta = 1.5$ simultaneously---high risk aversion with moderate intertemporal substitution.

\begin{calloutnote}[The Unit Elasticity Benchmark]
When $\gamma = \theta = 1$, both aggregations become geometric means (Cobb-Douglas), and preferences reduce to
\[
V_t = c_t^{1-\beta} \cdot \exp\left( \beta \, \E_t[\log V_{t+1}] \right).
\]
Maximizing $\log V_t$ gives the additively separable objective $\E_t[\sum_s \beta^{s-t} \log c_s]$. This is the log-utility case from Section~\ref{subsec:log-utility}, which admits closed-form portfolio solutions.
\end{calloutnote}

% --------------------------------------------------------------
\subsection{Preference for Early or Late Resolution of Uncertainty}
\label{subsec:resolution-timing}
% --------------------------------------------------------------

A distinctive feature of Epstein-Zin preferences is that they encode attitudes toward the \emph{timing} of uncertainty resolution, not just its presence.

\paragraph{The Key Parameter.}

Define
\begin{equation}
\psi \equiv \frac{1 - \gamma}{1 - 1/\theta} = \frac{\rho_{\text{risk}}}{\rho_{\text{time}}}.
\label{eq:psi-definition}
\end{equation}
This ratio determines preferences over the timing of resolution:
\begin{itemize}[nosep]
    \item $\psi < 1$ (equivalently, $\gamma > 1/\theta$): \textbf{Preference for early resolution}. The agent prefers to learn outcomes sooner rather than later, even if no action can be taken.
    \item $\psi > 1$ (equivalently, $\gamma < 1/\theta$): \textbf{Preference for late resolution}. The agent prefers to delay learning outcomes.
    \item $\psi = 1$ (equivalently, $\gamma = 1/\theta$): \textbf{Indifference}. This is the time-separable case---timing of resolution is irrelevant.
\end{itemize}

\paragraph{Intuition.}

Consider two lotteries that deliver the same consumption streams with the same probabilities, but differ in when uncertainty is resolved:
\begin{itemize}[nosep]
    \item \textbf{Early resolution}: A coin is flipped today; you learn tomorrow's consumption now.
    \item \textbf{Late resolution}: The same coin flip occurs, but you learn the outcome tomorrow.
\end{itemize}

With time-separable utility, these are equivalent---only the probability distribution over final outcomes matters. With Epstein-Zin preferences and $\gamma > 1/\theta$, early resolution is preferred. The agent values ``knowing where they stand'' even before acting.

\begin{calloutimportant}[Why Timing Matters for Portfolio Choice]
Preference for early resolution ($\gamma > 1/\theta$) has important implications:
\begin{itemize}[nosep]
    \item Agents value assets that resolve uncertainty quickly (e.g., options, insurance).
    \item Long-horizon investors may behave differently than short-horizon investors \emph{even with identical terminal objectives}.
    \item The equity premium reflects not just risk aversion but also compensation for bearing persistent uncertainty.
\end{itemize}
This is one reason Epstein-Zin preferences have become central to modern asset pricing.
\end{calloutimportant}

% --------------------------------------------------------------
\subsection{The Ordinally Equivalent Formulation}
\label{subsec:ordinal-ez}
% --------------------------------------------------------------

For analytical work, it is often convenient to use an ordinally equivalent transformation of \eqref{eq:epstein-zin}.

\paragraph{The Transformed Value Function.}

Define $\tilde{V}_t = V_t^{1-\gamma}$ when $\gamma \neq 1$. The recursion becomes
\begin{equation}
\tilde{V}_t = \left[ (1-\beta) c_t^{1-1/\theta} + \beta \left( \E_t[\tilde{V}_{t+1}] \right)^{\frac{1-1/\theta}{1-\gamma}} \right]^{\frac{1-\gamma}{1-1/\theta}}.
\label{eq:ez-transformed}
\end{equation}

\paragraph{The Weil Parametrization.}

An alternative form, due to Weil (1990), writes
\begin{equation}
U_t = \left[ (1-\beta) c_t^{\rho} + \beta \left( \E_t[U_{t+1}^{\alpha}] \right)^{\rho/\alpha} \right]^{1/\rho},
\label{eq:weil-form}
\end{equation}
where $\rho = 1 - 1/\theta$ and $\alpha = (1-\gamma)/(1-1/\theta)$. This makes the nested CES structure explicit: power $\rho$ for intertemporal aggregation, power $\alpha \cdot \rho$ effectively for risk aggregation.

\paragraph{Special Cases.}

\begin{center}
\renewcommand{\arraystretch}{1.4}
\begin{tabular}{@{}lccl@{}}
\toprule
\textbf{Case} & $\gamma$ & $\theta$ & \textbf{Utility Form} \\
\midrule
Log utility & $1$ & $1$ & $\E[\sum_t \beta^t \log c_t]$ \\
Time-separable CRRA & $\gamma$ & $1/\gamma$ & $\E[\sum_t \beta^t c_t^{1-\gamma}/(1-\gamma)]$ \\
Risk-neutral, unit IES & $0$ & $1$ & $\E[\sum_t \beta^t c_t] / (1-\beta)$ \\
High RA, high IES & $10$ & $1.5$ & Full Epstein-Zin recursion \\
\bottomrule
\end{tabular}
\end{center}

